\scanpage{13}
Dans ce cas favorable
\diagram{\begin{tikzcd}[ampersand replacement=\&,column sep=0.2em,row sep=1.4em]
\displaystyle\sum_i(-1)^i\frac{i+\rho}{2}b_i=\frac{\rho}{2}\chi(E)+\frac12 \& \displaystyle\underbrace{\sum_{i=0}^{2n}(-1)^i i b_i} \arrow[d,no head] \\
\& \displaystyle\sum_{i=0}^{n-1}\left[(-1)^i i b_i+(-1)^{2n-i}(2n-i)b_{2n-i}\right]+(-1)^n n b_n
\end{tikzcd}}
comme $b_i=b_{2n-i}$
\[
2n\sum_{i=0}^{n-1}(-1)^i b_i+(-1)^n n b_n=n\sum_{i=0}^{2n}(-1)^i b_i=n\chi(E)
\]
\[
\sum_i(-1)^i\frac{i+\rho}{2}b_i=\deleted{\unsure{n+\rho}}\frac{n+\rho}{2}\chi(E)
\]
donc
\[
\delta(E)=\unsure{\varepsilon(E)|\delta(E)}p^{\frac{n+\rho}{2}\chi(E)}
\qquad\varepsilon(E)=\pm1
\]
Par suite l'équation fonctionnelle devient dans ce cas
\[
(\unsure{2|1})\qquad L_{DE}(t)=\varepsilon(E)\left(-p^{\frac{n+\rho}{2}}\frac1t\right)^{\chi(E)}L_E(t^{-1}).
\]
Ou encore, si on admet que
\[
\check E\deleted{\unsure{DE}}\simeq E(\rho)
\]
globalement sur $X$ \deleted{\unsure{alors}}
\inserted{[il suffit que ce soit vrai sur le corps des fonctions de $X$, \unsure{donc c'est bon, donc tjrs vrai}]}
\[
\deleted{\unsure{L_E}}D(E)\simeq\deleted{\unsure{\check E\otimes\mathbb Q_\ell(n)}}\check E(n)[2n]\simeq E(\rho+n)[2n].
\]
\margin{\diagram{\begin{tikzcd}[ampersand replacement=\&,column sep=6em,row sep=0.4em]
{[2n]}\arrow[d,no head] \& {[2n]}\arrow[d,no head]\\
\text{shift} \& \text{shift}
\end{tikzcd}}}
et par suite
\[
\deleted{L_{DE}\simeq L_{E(\rho+n)}=L_{\mathbf R f_*(E(\rho+n))}}
\]
\[
\mathbf R f_*(D(E))\simeq\mathbf R f_*(E)(\rho+\deleted{2}n)[2n]
\]
d'où
\[
L_{DE}(t)=L_{\mathbf R f_*D(E)}(t)\simeq L_{\mathbf R f_*(E)}\left(p^{-(\rho+n)}t\right).
\]
\deleted{et il vient}
\mspage{12}
